\documentclass[11pt]{article}

\usepackage[a4paper,margin=1in]{geometry}
\usepackage{amsmath,amssymb,mathtools}
\usepackage[hidelinks]{hyperref}

\title{Selected Boundary-Layer Theorem}
\author{}
\date{}

\begin{document}
\maketitle

This note implements the hybrid route suggested by the Plan 1 and Plan 2 agent
reports. The guiding point is that the pure level-set method loses control at
the boundary for arbitrary rough domains, while the Plan 1 selected minimizers
have exactly the regularity needed to take small positive torsion levels as
controlled parallel surfaces.

The result below is split into two parts: an unconditional near-boundary
good-level theorem, using only the exact level-set identity and the profile-gap
estimate; and a selected-minimizer upgrade whose Serrin input is reduced in
\texttt{weighted-serrin-collar-reduction.md}.

Throughout,
\[
-\Delta u=1\quad\hbox{in }U,\qquad u=0\quad\hbox{on }\partial U,
\]
\[
E_t=\{u>t\},\qquad m(t)=|E_t|,\qquad
L=|U|,\qquad \delta_T(U)=E(U)-E(B_L),
\]
where \(B_L\) is the ball of volume \(L\). Write
\[
c_n=n^2\omega_n^{2/n}.
\]
For regular levels define
\[
D_H(t)=
\left(\int_{\{u=t\}}\frac1{|\nabla u|}\right)
\left(\int_{\{u=t\}}|\nabla u|\right)
-P(E_t)^2
\]
and
\[
D_I(t)=P(E_t)^2-c_n m(t)^{2-2/n}.
\]
The exact identity from \texttt{level-set-deficit-identity.md} is
\[
2\delta_T(U)
=
\int_0^{\|u\|_\infty}
\frac{D_H(t)+D_I(t)}
{c_n m(t)^{1-2/n}}\,dt.
\]

\section{Near-Boundary Good Levels}

Let
\[
R_L=\left(\frac{L}{\omega_n}\right)^{1/n},
\qquad
b(s)=u_{B_L}^*(s)
=\frac{R_L^2-(s/\omega_n)^{2/n}}{2n}.
\]
For \(0<\eta<L/4\), set
\[
S_\eta=[L-2\eta,L-\eta],
\qquad
T_\eta=u^*(S_\eta).
\]
Equivalently, \(t\in T_\eta\) means
\[
L-2\eta\le m(t)\le L-\eta.
\]

\paragraph{Proposition 1.1: boundary-slab good level.}
There are dimensional constants \(c,C>0\) such that if
\[
0<\eta<L/4,\qquad
\delta_T(U)\le c R_L^2\eta,
\]
then \(T_\eta\) contains a regular value \(t_\eta\) satisfying
\[
\eta\le |U\setminus E_{t_\eta}|\le 2\eta
\]
and
\[
D_H(t_\eta)+D_I(t_\eta)
\le
C\,c_n\,L^{1-2/n}\frac{\delta_T(U)}{R_L^2\eta/L}.
\]
Equivalently, since \(R_L^2\simeq_n L^{2/n}\),
\[
D_H(t_\eta)+D_I(t_\eta)
\le
C_n\frac{L^{2-4/n}}{\eta}\,\delta_T(U).
\]
Moreover the chosen level lies at the correct boundary height:
\[
c_n\frac{R_L^2}{L}\eta
\le t_\eta\le
C_n\frac{R_L^2}{L}\eta .
\]

\paragraph{Proof.}
Let \(v(s)=u^*(s)\) and \(h(s)=b(s)-v(s)\). The profile-gap estimate gives
\[
0\le h(s)\le \frac{2\delta_T(U)}{s}.
\]
On \(S_\eta\), \(s\ge L/2\), hence
\[
0\le h(s)\le \frac{4\delta_T(U)}{L}.
\]
For \(s\in S_\eta\), elementary bounds for \((1-\rho)^{2/n}\),
\(0\le \rho\le 1/2\), give
\[
c_n\frac{R_L^2}{L}\eta
\le b(s)\le
C_n\frac{R_L^2}{L}\eta.
\]
After decreasing \(c\), the hypothesis
\(\delta_T(U)\le c R_L^2\eta\) makes \(v(s)=b(s)-h(s)\) comparable to
\(R_L^2\eta/L\) throughout the slab.

The same estimate gives a lower bound for the length of \(T_\eta\). Indeed,
\[
|T_\eta|
=v(L-2\eta)-v(L-\eta)
\]
and
\[
b(L-2\eta)-b(L-\eta)\ge c_n\frac{R_L^2}{L}\eta.
\]
The oscillation of \(h\) across the slab is at most \(4\delta_T(U)/L\), so,
again by the smallness assumption,
\[
|T_\eta|\ge c_n\frac{R_L^2}{L}\eta.
\]
Now average the exact identity over \(T_\eta\). Since
\(m(t)\in[L-2\eta,L-\eta]\subset[L/2,L]\) on \(T_\eta\), there exists a
regular \(t_\eta\in T_\eta\) with
\[
D_H(t_\eta)+D_I(t_\eta)
\le
\frac{2\delta_T(U)}
{\int_{T_\eta}\frac{dt}{c_n m(t)^{1-2/n}}}
\le
C_n\frac{L^{2-4/n}}{\eta}\delta_T(U).
\]
The volume-layer bound follows directly from the definition of \(T_\eta\).

\section{Transfer Back to the Original Set}

The selected level is useful only if stability of \(E_{t_\eta}\) can be
returned to \(U\).

\paragraph{Lemma 2.1: boundary-layer asymmetry transfer.}
For every \(t>0\), with \(\eta=L-m(t)\),
\[
\mathcal A(U)\le \mathcal A(E_t)+2\frac{\eta}{L}.
\]

\paragraph{Proof.}
Let \(B_s\) be an optimal ball for \(E_t\), where \(s=m(t)=L-\eta\), and let
\(B_L\) be the concentric dilation of \(B_s\) with volume \(L\). Since
\(E_t\subset U\),
\[
|U\Delta B_L|
\le |U\setminus E_t|+|E_t\Delta B_s|+|B_L\setminus B_s|.
\]
Divide by \(L\). The first and third terms are both \(\eta/L\), while the
middle term is \((s/L)\mathcal A(E_t)\le \mathcal A(E_t)\).

\paragraph{Corollary 2.2: good-level transfer estimate.}
If \(t_\eta\) is the level from Proposition 1.1, then
\[
\mathcal A(U)^2
\le
2\mathcal A(E_{t_\eta})^2
+8\frac{\eta^2}{L^2}.
\]
By the quantitative isoperimetric inequality applied to \(E_{t_\eta}\),
\[
\mathcal A(E_{t_\eta})^2
\le
C_n\frac{D_I(t_\eta)}{m(t_\eta)^{2-2/n}},
\]
and therefore
\[
\mathcal A(U)^2
\le
C_n\frac{D_I(t_\eta)}{L^{2-2/n}}
+8\frac{\eta^2}{L^2}.
\]
Choosing
\[
\eta=L\,\delta_{\rm rel}^{1/2},
\qquad
\delta_{\rm rel}:=\frac{\delta_T(U)}{L^{1+2/n}},
\]
one obtains the scale-invariant estimate
\[
\mathcal A(U)^2
\le
C_n\delta_{\rm rel}^{1/2}
\]
from level-set information alone. This is the known boundary-layer loss: it
gives the wrong exponent. The selected-minimizer regularity is needed to
replace this average estimate by a boundary trace estimate.

\section{Selected-Minimizer Boundary Trace}

Assume now that \(U=\{u>0\}\) is a Plan 1 selected minimizer in the graph
regime. Thus \(\partial U\) is \(C^{1,\gamma}\), the Bernoulli coefficient
\[
q(y)=|\nabla u(y)|,\qquad y\in\partial U,
\]
satisfies
\[
0<q_-\le q\le q_+,\qquad [q]_{C^\gamma(\partial U)}\le M,
\]
and the interior normal exponential map is a \(C^{1,\gamma}\) diffeomorphism in
a fixed collar of \(\partial U\).

\paragraph{Lemma 3.1: positive levels are controlled parallel surfaces.}
There are \(t_0>0\) and \(C\), depending only on the selected-minimizer
regularity constants, such that for \(0<t<t_0\)
\[
\{u=t\}
=
\left\{
y-\frac{t}{q(y)}\nu_U(y)+R_t(y):
y\in\partial U
\right\},
\]
with
\[
|R_t(y)|\le C t^{1+\gamma}.
\]
Consequently,
\[
|U\setminus E_t|
=
t\int_{\partial U}\frac1q\,d\mathcal H^{n-1}
+O(t^{1+\gamma}),
\]
\[
P(E_t)=P(U)+O(t),
\]
and
\[
\bigl|\mathcal A(E_t)-\mathcal A(U)\bigr|\le C t.
\]

\paragraph{Proof.}
In the collar, write points as \(x=y-r\nu_U(y)\). Boundary \(C^{1,\gamma}\)
regularity and the Bernoulli condition give
\[
u(y-r\nu_U(y))=q(y)r+O(r^{1+\gamma})
\]
uniformly in \(y\). Since \(q\ge q_->0\), the equation \(u=t\) has a unique
solution
\[
r_t(y)=\frac{t}{q(y)}+O(t^{1+\gamma})
\]
by the elementary inverse estimate for \(C^{1,\gamma}\) perturbations of a
linear function. The area and volume expansions follow from the area formula
for the normal map. The asymmetry estimate follows from
\(|E_t\Delta U|=|U\setminus E_t|=O(t)\) and the fixed-volume Lipschitz
continuity of Fraenkel asymmetry under symmetric difference, after the harmless
volume normalization.

\section{The Conditional Sharp Theorem}

The preceding lemmas isolate the only missing analytic input.

\paragraph{Input S: weighted Serrin stability on selected collars.}
For selected minimizers in the graph regime, assume the following stability
estimate holds uniformly for all sufficiently small regular levels:
\[
\mathcal A(E_t)^2
\le
C\left(
\frac{D_H(t)+D_I(t)}{L^{2-2/n}}
+\frac{|U\setminus E_t|^2}{L^2}
\right).
\]
The \(D_I\) part is the quantitative isoperimetric inequality. The real content
is to upgrade the weighted variance identity
\[
\int_{\{u=t\}}
\frac{(|\nabla u|-\bar f_t)^2}{|\nabla u|}
=
\frac{m(t)}{P(E_t)^2}D_H(t)
\]
to the boundary norm required by a Serrin-stability theorem, using the uniform
\(C^\gamma\) estimates supplied by Plan 1.

The follow-up note \texttt{weighted-serrin-collar-reduction.md} observes that
this specific Input S is already implied by the quantitative isoperimetric
inequality through \(D_I(t)\). The useful role of \(D_H(t)\) is diagnostic: it
provides approximate Serrin data, but it is not the sharp Fraenkel bottleneck in
the selected graph regime.

\paragraph{Theorem 4.1: selected boundary-layer theorem, conditional on Input S.}
Let \(U\) be a Plan 1 selected minimizer in the graph regime, with
\(|U|=L\), and let
\[
\delta_{\rm rel}:=\frac{\delta_T(U)}{L^{1+2/n}}.
\]
Assume Input S. If \(\delta_{\rm rel}\) is sufficiently small, then there is a
regular level \(t\) such that
\[
|U\setminus E_t|\le C L\delta_{\rm rel}^{1/2},
\]
\[
E_t\hbox{ is a regular approximate Serrin domain},
\]
and
\[
\mathcal A(U)^2
\le
C\left(
\frac{D_H(t)+D_I(t)}{L^{2-2/n}}
+\delta_{\rm rel}
\right).
\]
In particular,
\[
\mathcal A(U)^2\le C\delta_{\rm rel}^{1/2}
\]
unconditionally from the averaged level-set bound, and the sharp estimate
\[
\mathcal A(U)^2\le C\delta_{\rm rel}
\]
would follow once the boundary deficit propagation step finds a boundary level
with
\[
D_H(t)+D_I(t)\le C L^{2-2/n}\delta_{\rm rel}.
\]

\paragraph{Proof.}
Choose \(\eta=L\delta_{\rm rel}^{1/2}\). Proposition 1.1 gives a regular
\(t\) with \(|U\setminus E_t|\le 2\eta\) and controlled level deficit. Lemma
3.1 ensures that, for small \(\delta_{\rm rel}\), this level is a regular
parallel surface in the selected collar. Input S gives the displayed estimate
for \(\mathcal A(E_t)^2\), and Lemma 2.1 transfers it back to \(U\). The
boundary-layer volume contributes exactly \(\eta^2/L^2=\delta_{\rm rel}\).

The last statement identifies the remaining sharpness gap: the current
averaging argument gives a \(1/\eta\) loss, while the selected-minimizer trace
theorem should use the small-collar regularity to avoid averaging over a slab
whose length is only \(O(\eta)\).

\section{Implementation Consequence for the Full Strategy}

This note makes the Plan 1/Plan 2 hybrid route precise:
\begin{enumerate}
\item Plan 1 selects \(U\) and preserves the quotient
\((E(\Omega)-E(B))/\alpha(\Omega)\).
\item Plan 1 graph entry gives a \(C^{1,\gamma}\) collar and Bernoulli bounds.
\item Plan 2 gives the exact integrated level-set identity.
\item Proposition 1.1 finds a genuine near-boundary level.
\item Lemma 3.1 turns that level into a controlled trace of \(\partial U\).
\item The remaining non-formal input is boundary deficit propagation:
\(D_I(0)+D_H(0)\lesssim E(U)-E(B)\), or an equivalent Plan 1 near-spherical
expansion.
\end{enumerate}
Thus the most promising next concrete task is the Plan 1 graph-entry and
near-spherical closure, not another selection lemma or a stronger Serrin theorem
for rough levels.

\end{document}
